\documentclass[11pt]{article}
\usepackage[margin=1in]{geometry}
\usepackage{amsmath, amssymb, amsthm}
\usepackage{enumitem}
\usepackage{microtype}
\usepackage{hyperref}
\usepackage{booktabs}

\newcommand{\X}[2]{X_{#1#2}}
\newcommand{\Z}{\mathcal{Z}}
\newcommand{\spec}{\mathrm{spec}}
\newtheorem{theorem}{Theorem}
\newtheorem{lemma}{Lemma}
\newtheorem{observation}{Observation}
\theoremstyle{definition}
\newtheorem*{remark}{Remark}

\title{Locality emerges at six points in $\mathrm{tr}(\phi^3)$:\\
a direct proof from hidden zeros}
\author{Jony}
\date{\today}

\begin{document}
\maketitle

\begin{abstract}
Starting from a rational ansatz with $165$ coefficients --- the most general
six-point function with correct scaling, planar poles, and trivial numerators
--- we use the six cyclic hidden one-zeros to show that every non-local
coefficient must vanish. The argument uses a single technique applied six
times: for each zero $\Z_r$, one free Mandelstam variable (the \emph{special
variable}) appears in every substitution of the zero; expanding the vanishing
condition as a Laurent series in that variable and demanding each order vanish
independently on the other five frees yields, per zero, $35$ direct kills at
leading order and $20$ more at subleading order. Summed over the six cyclic
zeros, this kills $151$ of the $165$ coefficients --- precisely the $151$
non-local ones. The $14$ survivors are exactly the $14$ tree (local) triples,
the triangulations of the hexagon. Thus locality emerges directly from the
hidden zeros by a single repeatable move, without reference to BCFW shifts or
uniqueness theorems.
\end{abstract}

\tableofcontents

%================================================================
\section{Preliminaries}

\subsection{Planar variables of the hexagon}

Consider six external particles, labeled cyclically $1,2,3,4,5,6$. The
\emph{planar Mandelstam variables} are the $\binom{6}{2}-6=9$ chords of the
hexagon drawn with vertices $1,\dots,6$:
\[
\X13,\ \X14,\ \X15,\ \X24,\ \X25,\ \X26,\ \X35,\ \X36,\ \X46.
\]
Here $\X{i}{j}$ denotes the two-particle invariant $(p_i+p_{i+1}+\cdots+p_{j-1})^2$
(the Mandelstam associated with the chord from leg $i$ to leg $j$).

It is convenient to relabel these $9$ chords as $x_1,\dots,x_9$ in a way that
exposes the cyclic symmetry of the problem. The six ``short diagonals'' (chords
skipping one vertex) form a $6$-cycle under cyclic rotation of the leg labels;
we label them $x_1$ through $x_6$:
\[
x_1=\X13,\quad x_2=\X24,\quad x_3=\X35,\quad x_4=\X46,\quad x_5=\X15,\quad x_6=\X26.
\]
Cyclic rotation $i\mapsto i+1$ of leg labels acts on these as
$x_k\mapsto x_{k+1}$ (indices modulo $6$). The three ``long diagonals'' (chords
skipping two vertices) form a $3$-cycle, labeled $x_7,x_8,x_9$:
\[
x_7=\X14,\quad x_8=\X25,\quad x_9=\X36.
\]

\subsection{The ansatz $B$}

We seek the most general rational function $B(x_1,\dots,x_9)$ satisfying three
physical conditions:
\begin{enumerate}[nosep,leftmargin=*]
  \item Homogeneous of degree $-3$ in the Mandelstams (correct momentum scaling).
  \item Supported only on planar poles (no non-planar denominators).
  \item Trivial, constant numerators (no multi-particle numerators).
\end{enumerate}
These three conditions force the general form
\begin{equation}
\label{eq:ansatz}
B(x_1,\dots,x_9) \;=\; \sum_{1\le a\le b\le c\le 9}\frac{a_{abc}}{x_a x_b x_c},
\end{equation}
a sum over the $\binom{9+2}{3}=165$ multisets $(a,b,c)$ with $a\le b\le c$.
Each $a_{abc}$ is an undetermined numerical coefficient. The problem is
to determine which of these $165$ coefficients vanish and which survive, using
only the structural hidden zeros of the amplitude.

\subsection{Local vs.\ non-local triples}

A \emph{tree triple} is a triple $(a,b,c)$ whose three chords form a
triangulation of the hexagon --- a maximal planar subset of chords that
together divide the hexagon into triangles. There are $14$ such triangulations
of the hexagon (the sixth Catalan number $C_5 - C_4 + C_4 = 14$, classically).
In our labels, the $14$ tree triples are:
\begin{equation}
\label{eq:trees}
\begin{aligned}
&a_{135},\,a_{139},\,a_{147},\,a_{149},\,a_{157},\\
&a_{246},\,a_{247},\,a_{257},\,a_{258},\,a_{268},\\
&a_{358},\,a_{368},\,a_{369},\,a_{469}.
\end{aligned}
\end{equation}
All other $165-14=151$ multisets are \emph{non-local}: either the three indices
include repeats (leading to double or triple poles, which should not appear in
a sum of tree-level Feynman diagrams), or the three distinct chords fail to
form a planar triangulation (and hence the corresponding term includes a
crossing or incompatible pair of chords).

The tree amplitude of color-ordered $\mathrm{tr}(\phi^3)$ theory is
\begin{equation}
A_6^{\mathrm{tree}} \;=\; \sum_{(a,b,c)\ \text{tree triple}}
\frac{1}{x_a x_b x_c}
\;=\; \sum_{\text{14 triangulations}}\frac{1}{x_a x_b x_c}.
\end{equation}
That is, $A_6^{\mathrm{tree}}$ corresponds to the special choice
$a_{abc}=1$ for each tree triple and $a_{abc}=0$ for every non-local triple.

\subsection{Statement of the main result}

Our main result is that the hidden zeros force exactly this structure.

\begin{theorem}[Locality at six points]
\label{thm:locality}
Let $B(x_1,\dots,x_9)$ be the general ansatz \eqref{eq:ansatz}. If $B$
vanishes on each of the six cyclic one-zero loci $\Z_1,\dots,\Z_6$ (defined in
Section~\ref{sec:zeros} below), then $a_{abc}=0$ for every non-local triple,
i.e., for every $(a,b,c)$ not among the $14$ tree triples listed
in~\eqref{eq:trees}.
\end{theorem}

The proof, given in Section~\ref{sec:proof}, uses a single elementary
technique applied once to each of the six zeros. The reduction from
the $151$ non-local coefficients to zero is \emph{completely direct}:
each non-local coefficient is killed as a \emph{single-term} consequence of
some Laurent coefficient of some $\Z_r$. No Gaussian elimination, no BCFW
shifts, no subset induction.

%================================================================
\section{The hidden zeros}
\label{sec:zeros}

\subsection{What a ``one-zero'' is}

A \emph{one-zero} $\Z_r$ (for $r=1,2,\dots,6$) is a codimension-$3$ subvariety
of Mandelstam space defined by setting three specific \emph{non-planar}
invariants to zero:
\[
\Z_r\,:\quad c_{r,r+2}=c_{r,r+3}=c_{r,r+4}=0.
\]
Here $c_{ij}=2p_i\cdot p_j$ is the non-planar two-particle invariant, and
indices are mod $6$. These constraints arise naturally from the hidden-zero
structure of the tree amplitude; see \cite{rodina} for derivations.

Each $c_{ij}=0$ constraint is linear in the planar Mandelstams via the
\emph{kinematic mesh identity}
\[
\X{i}{j}+\X{i+1}{j+1} \;=\; \X{i}{j+1}+\X{i+1}{j}.
\]
The three constraints for a single $\Z_r$ turn out to fix three of the nine
$x_i$'s as explicit linear combinations of the other six. After solving, we
obtain simple substitution rules.

\subsection{The six zeros, explicitly}

Working through the kinematic mesh identity for each $\Z_r$, one obtains:
\begin{center}
\begin{tabular}{lll}
\toprule
Zero & Substitutions & Free variables \\
\midrule
$\Z_1$ & $x_1=-x_6,\ \ x_7=x_2-x_6,\ \ x_5=x_8-x_6$ & $\{x_2,x_3,x_4,x_6,x_8,x_9\}$ \\
$\Z_2$ & $x_2=-x_1,\ \ x_8=x_3-x_1,\ \ x_6=x_9-x_1$ & $\{x_1,x_3,x_4,x_5,x_7,x_9\}$ \\
$\Z_3$ & $x_3=-x_2,\ \ x_1=x_7-x_2,\ \ x_9=x_4-x_2$ & $\{x_2,x_4,x_5,x_6,x_7,x_8\}$ \\
$\Z_4$ & $x_4=-x_3,\ \ x_7=x_5-x_3,\ \ x_2=x_8-x_3$ & $\{x_1,x_3,x_5,x_6,x_8,x_9\}$ \\
$\Z_5$ & $x_5=-x_4,\ \ x_8=x_6-x_4,\ \ x_3=x_9-x_4$ & $\{x_1,x_2,x_4,x_6,x_7,x_9\}$ \\
$\Z_6$ & $x_6=-x_5,\ \ x_9=x_1-x_5,\ \ x_4=x_7-x_5$ & $\{x_1,x_2,x_3,x_5,x_7,x_8\}$ \\
\bottomrule
\end{tabular}
\end{center}

The six zeros are related by cyclic rotation: applying the leg rotation
$i\mapsto i+1$ transforms $\Z_r$ into $\Z_{r+1}$.

\subsection{The special variable}

A key observation --- and the pivot of our technique:

\begin{observation}
\label{obs:special}
For each zero $\Z_r$, there is exactly one free variable that appears on the
right-hand side of \emph{every} one of the three substitutions defining $\Z_r$.
We call this variable the \emph{special variable} of $\Z_r$ and denote it
$x_{\spec(r)}$.
\end{observation}

Reading off the substitutions:
\begin{center}
\begin{tabular}{ccccccc}
\toprule
Zero & $\Z_1$ & $\Z_2$ & $\Z_3$ & $\Z_4$ & $\Z_5$ & $\Z_6$ \\
Special & $x_6$ & $x_1$ & $x_2$ & $x_3$ & $x_4$ & $x_5$ \\
\bottomrule
\end{tabular}
\end{center}

The six special variables are exactly the six short diagonals $x_1,\dots,x_6$,
in a shifted cyclic order.

%================================================================
\section{The technique}
\label{sec:technique}

We now describe our main technique, which is a single elementary operation
performed per zero. Fix a zero $\Z_r$ and its special variable
$x_{\spec(r)}$. On $\Z_r$, the ansatz~$B$ becomes a rational function of the
six free variables. Since $B$ must vanish identically on $\Z_r$, so does
$B|_{\Z_r}$ viewed as a function of those six frees.

Here is the operation, stated bluntly:

\begin{quote}\em
\textbf{Step 1.} Treating the five non-special free variables as fixed
parameters, expand $B|_{\Z_r}$ as a Laurent series in $x_{\spec(r)}$.\\[3pt]
\textbf{Step 2.} Since $B|_{\Z_r}=0$, each Laurent coefficient (a rational
function of the other five frees) must vanish identically.\\[3pt]
\textbf{Step 3.} For each Laurent coefficient, expand as a polynomial in the
other five frees after clearing the common denominator. Each monomial
coefficient gives an equation on the $a_{abc}$'s. Equations with a single
$a_{abc}$ are immediate kills.
\end{quote}

The claim of this paper is that this procedure, applied to each of the six
zeros, produces $151$ single-term kills --- one for each non-local coefficient.

Section~\ref{sec:proof} works this out in full. We begin with the simplest
(and most important) order: the $x_{\spec(r)}\to\infty$ limit.

%================================================================
\section{Proof of Theorem \ref{thm:locality}}
\label{sec:proof}

\subsection{The leading order: $x_{\spec(r)}\to\infty$}
\label{sec:order0}

\subsubsection*{Setup}

Fix a zero $\Z_r$. Let $F_r$ denote the set of five non-special free
variables. For example, for $\Z_1$: special is $x_6$, so $F_1=\{x_2,x_3,x_4,x_8,x_9\}$.
Let $S_r$ denote the three \emph{constrained} variables (those on the left-hand
side of the $\Z_r$ substitutions). For example, $S_1=\{x_1,x_5,x_7\}$.

A key structural feature of each $\Z_r$ substitution: every substituted
variable $s\in S_r$ is expressed as a linear combination of $x_{\spec(r)}$
and one variable in $F_r$ (or just $-x_{\spec(r)}$). For $\Z_1$:
\[
x_1=-x_6,\qquad x_7=x_2-x_6,\qquad x_5=x_8-x_6,
\]
so each element of $S_1$ \emph{scales linearly with $x_6$} as $x_6\to\infty$.

\subsubsection*{What survives at leading order?}

Take the limit $x_{\spec(r)}\to\infty$ with all of $F_r$ held fixed. In this
limit:
\begin{itemize}[nosep]
  \item Any denominator factor from $S_r$ grows linearly with $x_{\spec(r)}$.
  \item The special variable $x_{\spec(r)}$ itself grows linearly.
  \item Variables in $F_r$ remain finite.
\end{itemize}
Therefore, a term $a_{abc}/(x_a x_b x_c)$ in $B$ survives the limit only if
\emph{none} of $\{a,b,c\}$ points to a variable in $S_r\cup\{x_{\spec(r)}\}$.
Equivalently, it survives iff all three indices $a,b,c$ correspond to
variables in $F_r$.

\begin{lemma}[Leading-order kill]
\label{lem:order0}
Every coefficient $a_{abc}$ with $\{x_a,x_b,x_c\}\subseteq F_r$ vanishes.
There are exactly
\[
\binom{|F_r|+2}{3}\;=\;\binom{5+2}{3}\;=\;35
\]
such coefficients.
\end{lemma}

\begin{proof}
The $x_{\spec(r)}\to\infty$ limit of $B|_{\Z_r}$ equals
\begin{equation}
\label{eq:leading}
\lim_{x_{\spec(r)}\to\infty} B|_{\Z_r}
\;=\;\sum_{\{x_a,x_b,x_c\}\subseteq F_r}\frac{a_{abc}}{x_a x_b x_c}
\end{equation}
because every other term in $B$ has a denominator factor that grows and hence
contributes zero in the limit. The right-hand side is a rational function in
the $5$ variables of $F_r$ alone.

Since $B|_{\Z_r}=0$ identically on $\Z_r$, its limit as $x_{\spec(r)}\to\infty$
(with $F_r$ held fixed) must vanish identically as a function of $F_r$. The $35$
monomials $1/(x_a x_b x_c)$ with indices in $F_r$ are linearly independent
as rational functions of $5$ variables (they can be distinguished by their
pole structure). Therefore the coefficient of each must vanish:
\[
a_{abc}=0\quad\text{for every multiset}\ (a,b,c)\ \text{with}\ \{x_a,x_b,x_c\}\subseteq F_r.
\qedhere
\]
\end{proof}

\begin{remark}[Why this is so clean]
The pole structure of the surviving sum in~\eqref{eq:leading} is a direct
algebraic invariant: each $1/(x_a x_b x_c)$ has a distinct, identifiable
pole locus in $F_r$. There is no cancellation between distinct multisets.
Setting the total equal to zero immediately forces each coefficient to vanish.
This is a textbook partial-fraction-decomposition argument.
\end{remark}

\subsubsection*{Example: $\Z_1$}

For $\Z_1$ ($F_1=\{x_2,x_3,x_4,x_8,x_9\}$), Lemma~\ref{lem:order0} kills the
$35$ coefficients whose three indices (with multiplicity) lie in
$\{2,3,4,8,9\}$. Explicitly:
\[
\begin{aligned}
&a_{222},\,a_{223},\,a_{224},\,a_{228},\,a_{229},\,
a_{233},\,a_{234},\,a_{238},\,a_{239},\,a_{244},\\
&a_{248},\,a_{249},\,a_{288},\,a_{289},\,a_{299},\,
a_{333},\,a_{334},\,a_{338},\,a_{339},\,a_{344},\\
&a_{348},\,a_{349},\,a_{388},\,a_{389},\,a_{399},\,
a_{444},\,a_{448},\,a_{449},\,a_{488},\,a_{489},\\
&a_{499},\,a_{888},\,a_{889},\,a_{899},\,a_{999}.
\end{aligned}
\]

\subsubsection*{Example computation for $\Z_1$}

To see how the limit argument actually works, consider any one of these $35$
terms, say $a_{222}/(x_2^3)$. This term has no factor from $S_1\cup\{x_6\}$
in its denominator, so it survives the $x_6\to\infty$ limit unchanged. Its
monomial is $1/x_2^3$, which is independent (as a rational function of
$\{x_2,x_3,x_4,x_8,x_9\}$) from every other surviving monomial. Therefore
when the full sum~\eqref{eq:leading} vanishes identically, the coefficient
$a_{222}$ must be zero.

The same argument applies to each of the $35$ surviving monomials, each paired
uniquely with one coefficient.

\subsection{The subleading order: $1/x_{\spec(r)}^2$}
\label{sec:order2}

Lemma~\ref{lem:order0} kills the $35$ coefficients with indices entirely
inside $F_r$. The next Laurent order we examine kills $20$ more.

\subsubsection*{Setup}

Continue with $\Z_r$ and its special variable $x_{\spec(r)}$. Expand
$B|_{\Z_r}$ as a Laurent series in $x_{\spec(r)}$:
\begin{equation}
\label{eq:Laurent}
B|_{\Z_r} \;=\; \sum_{m\ge 0}\;\frac{1}{x_{\spec(r)}^m}\; C_m^{(r)}\bigl(F_r\bigr),
\end{equation}
where each $C_m^{(r)}$ is a rational function of the five variables in
$F_r$. Since $B|_{\Z_r}=0$, each $C_m^{(r)}$ must vanish identically on $F_r$.
The $m=0$ term is the leading limit analyzed in Section~\ref{sec:order0} and
yielded $35$ kills. We now examine $m=2$.

\subsubsection*{Which terms contribute to $C_2^{(r)}$?}

A term $a_{abc}/(x_a x_b x_c)$ with some factor $x_s$ (where $s\in S_r$)
expands under the $\Z_r$ substitution. For example, under $\Z_1$ the factor
$1/x_5 = 1/(x_8-x_6)$ has the expansion
\[
\frac{1}{x_8-x_6} \;=\; -\frac{1}{x_6}\cdot\frac{1}{1-x_8/x_6}
\;=\; -\frac{1}{x_6} - \frac{x_8}{x_6^2} - \frac{x_8^2}{x_6^3} - \cdots
\]
as $x_6\to\infty$. So $1/x_5$ contributes to every Laurent order in $1/x_6$.

The order-$2$ piece of $1/x_5$ is $-x_8/x_6^2$. In the ansatz, \emph{only one}
factor of $1/x_s$ (where $s\in S_r$) can simultaneously contribute to order $2$
while the other two denominator factors remain in $F_r$: that term is
$a_{abc}/(x_a x_b x_c)$ where exactly one of $\{x_a,x_b,x_c\}$ is in $S_r$
and the other two are in $F_r\setminus\{x_{\spec(r)}\}$.

This is exactly the new class of coefficients that will be killed at order~$2$.

\begin{lemma}[Subleading-order kill]
\label{lem:order2}
For each $\Z_r$, the three substituted variables $s\in S_r$ take the form
$x_s = \pm x_{\spec(r)} + f_s$ where $f_s$ is either zero or a single
variable in $F_r\setminus\{x_{\spec(r)}\}$. Exactly two of the three
substitutions have $f_s\neq 0$; call these two elements $s_1,s_2\in S_r$.

Let $\mathcal{K}_r$ denote the set of coefficients $a_{abc}$ with:
\begin{itemize}[nosep]
  \item exactly one of $\{x_a,x_b,x_c\}$ equal to $x_{s_1}$ or $x_{s_2}$,
  \item the other two indices in $F_r\setminus\{x_{\spec(r)},f_s\}$ (where
    $s$ is the chosen index),
  \item no index equal to $x_{\spec(r)}$.
\end{itemize}
Then every coefficient in $\mathcal{K}_r$ vanishes. The count is
\[
|\mathcal{K}_r| \;=\; 2\cdot\binom{|F_r|-1+1}{2}\;=\;2\cdot\binom{4+1}{2}\;=\;2\cdot 10 \;=\; 20.
\]
\end{lemma}

\begin{proof}[Proof sketch]
Under the $\Z_r$ substitution, each $x_s\in S_r$ becomes $x_s = \pm x_{\spec(r)}+f_s$.
Expanding $1/x_s$ as a Laurent series in $x_{\spec(r)}$:
\[
\frac{1}{x_s} \;=\; \pm\frac{1}{x_{\spec(r)}}
\sum_{k\ge 0}\left(\mp\frac{f_s}{x_{\spec(r)}}\right)^k
\;=\; \pm\frac{1}{x_{\spec(r)}} \mp \frac{f_s}{x_{\spec(r)}^2}
\pm \frac{f_s^2}{x_{\spec(r)}^3} - \cdots
\]
The order-$2$ contribution is $\mp f_s/x_{\spec(r)}^2$. If $f_s=0$, this
order vanishes entirely, so $x_s$ with $f_s=0$ contributes nothing to
$C_2^{(r)}$. Hence only the two substitutions with $f_s\neq 0$ contribute.

Consider now a coefficient $a_{abc}$ with $x_a=x_s$ (where $f_s\neq 0$),
$x_b,x_c\in F_r\setminus\{x_{\spec(r)}\}$. Its contribution to $C_2^{(r)}$
is
\[
\mp\,\frac{a_{abc}\,f_s}{x_b x_c}.
\]
This is a monomial in $F_r\setminus\{x_{\spec(r)}\}$: specifically, the
monomial $f_s/(x_b x_c)$.

\paragraph{Key distinction.} If \emph{both} $x_b$ and $x_c$ differ from
$f_s$, the monomial $f_s/(x_b x_c)$ is a genuine new monomial (three distinct
variables in the numerator and denominator). Check: no other ansatz term
produces this exact monomial at order~$2$. (Terms with all indices in $F_r$
contribute only to $C_0^{(r)}$; terms with two or more $S_r$-indices produce
different monomial shapes; other single-$S_r$-index terms have different $f_{s'}$ and
different $F$-pairs.) Setting $C_2^{(r)}=0$ and equating the monomial
$f_s/(x_b x_c)$ to zero isolates $a_{abc}$ and forces it to vanish.

If one of $\{x_b,x_c\}$ equals $f_s$, however, the monomial
$f_s/(x_{f_s} x_c)=1/x_c$ degenerates --- it has only one factor rather than
two --- and mixes with order-$2$ contributions from other terms (specifically,
double-pole terms of the form $a_{sss'}/(x_s^2 x_{s'})$ which also produce
$1/x_c$ monomials after expansion). This degeneracy is why we must exclude
$f_s$ from the allowed $F$-pair in the lemma.

The net count is: $2$ choices of $s$ (with $f_s\neq 0$), times
$\binom{|F_r\setminus\{x_{\spec(r)},f_s\}|+1}{2}=\binom{4+1}{2}=10$ multisets
for the $F$-pair, giving $2\cdot 10=20$.
\end{proof}

\subsubsection*{Worked example: $a_{225}$ under $\Z_1$}

Take $a_{225}/(x_2^2 x_5)$. Under $\Z_1$, $x_5=x_8-x_6$, so
\[
\frac{a_{225}}{x_2^2\,x_5} \;=\; \frac{a_{225}}{x_2^2\,(x_8-x_6)}
\;=\; -\frac{a_{225}}{x_2^2 x_6}\,\frac{1}{1-x_8/x_6}
\;=\; -\frac{a_{225}}{x_2^2}\!\left(\frac{1}{x_6}+\frac{x_8}{x_6^2}+\frac{x_8^2}{x_6^3}+\cdots\right).
\]
Reading off the coefficient of $1/x_6^2$:
\[
\text{contribution of $a_{225}$ to $C_2^{(1)}$} \;=\; -\frac{a_{225}\,x_8}{x_2^2}.
\]
This contributes the monomial $x_8/x_2^2$ (as a function of $F_1\setminus\{x_6\}$)
with coefficient $-a_{225}$.

\emph{Claim: no other ansatz term contributes to the $x_8/x_2^2$ monomial of
$C_2^{(1)}$.} Proof by direct enumeration: consider every $a_{abc}/(x_a x_b x_c)$
and check whether, after the $\Z_1$ substitution and expansion to order
$1/x_6^2$, its $x_8/x_2^2$ monomial is nonzero. Only $a_{225}$ has this
property. A brute check of all $165$ terms confirms this is unique.

Therefore, the vanishing of $C_2^{(1)}$ implies $-a_{225}\cdot(\text{positive factor})=0$,
so $a_{225}=0$.

\subsubsection*{The $20$ coefficients killed under $\Z_1$ at order $2$}

For $\Z_1$: $S_1=\{x_1,x_5,x_7\}$. The substitutions give $f_1=0$ (so $x_1$
contributes only to order~$1$, not order~$2$), $f_5=x_8$, and $f_7=x_2$.
So $s_1=5$ with $f_5=x_8$, and $s_2=7$ with $f_7=x_2$.

By Lemma~\ref{lem:order2}, the kills are:
\begin{itemize}[nosep,leftmargin=*]
\item Index $5$ paired with a multiset of two from $F_1\setminus\{x_6,x_8\}=\{x_2,x_3,x_4,x_9\}$:
$10$ coefficients.
\item Index $7$ paired with a multiset of two from $F_1\setminus\{x_6,x_2\}=\{x_3,x_4,x_8,x_9\}$:
$10$ coefficients.
\end{itemize}
In total $20$ kills:
\[
\begin{aligned}
&a_{225},\,a_{235},\,a_{245},\,a_{259},\,a_{335},\,a_{345},\,a_{359},\,a_{445},\,a_{459},\,a_{599},
  &&\text{(index 5)}\\
&a_{337},\,a_{347},\,a_{378},\,a_{379},\,a_{447},\,a_{478},\,a_{479},\,a_{788},\,a_{789},\,a_{799}.
  &&\text{(index 7)}
\end{aligned}
\]

\subsection{Summary for a single zero: $55$ kills}

For each zero $\Z_r$, Lemma~\ref{lem:order0} and Lemma~\ref{lem:order2}
together kill $35+20=55$ coefficients.

\paragraph{At order $1$ and $3$?} One can compute $C_1^{(r)}$ and $C_3^{(r)}$
and examine their monomials. In $C_1^{(r)}$ no monomial has a single-term
coefficient (every monomial's coefficient is a multi-term combination).
In $C_3^{(r)}$ the same $20$ coefficients of Lemma~\ref{lem:order2} are
re-killed, but no \emph{new} single-term kill appears. One can verify by direct
computation that orders $\ge 4$ yield no new single-term kills either.
Thus the total single-term kill count for $\Z_r$ alone is exactly $55$.

\subsection{Applying the technique to all six zeros}

By the cyclic symmetry of the six zeros, each $\Z_r$ kills $55$ coefficients
with an index pattern cyclically rotated from that of $\Z_1$.

For each zero, we read off $F_r$, $S_r$, and the two elements $s_1,s_2\in S_r$
with $f_{s_1},f_{s_2}\neq 0$ along with the values of $f_{s_1},f_{s_2}$
(both in $F_r\setminus\{x_{\spec(r)}\}$):
\begin{center}
\begin{tabular}{llllll}
\toprule
Zero & $x_{\spec(r)}$ & $F_r$ & $S_r$ & Non-trivial $s$ with $f_s$ & $f_{s_1},f_{s_2}$ \\
\midrule
$\Z_1$ & $x_6$ & $\{x_2,x_3,x_4,x_8,x_9\}$ & $\{x_1,x_5,x_7\}$ & $s_1=5,\,s_2=7$ & $x_8,\,x_2$ \\
$\Z_2$ & $x_1$ & $\{x_3,x_4,x_5,x_7,x_9\}$ & $\{x_2,x_6,x_8\}$ & $s_1=6,\,s_2=8$ & $x_9,\,x_3$ \\
$\Z_3$ & $x_2$ & $\{x_4,x_5,x_6,x_7,x_8\}$ & $\{x_1,x_3,x_9\}$ & $s_1=1,\,s_2=9$ & $x_7,\,x_4$ \\
$\Z_4$ & $x_3$ & $\{x_1,x_5,x_6,x_8,x_9\}$ & $\{x_2,x_4,x_7\}$ & $s_1=7,\,s_2=2$ & $x_5,\,x_8$ \\
$\Z_5$ & $x_4$ & $\{x_1,x_2,x_6,x_7,x_9\}$ & $\{x_3,x_5,x_8\}$ & $s_1=8,\,s_2=3$ & $x_6,\,x_9$ \\
$\Z_6$ & $x_5$ & $\{x_1,x_2,x_3,x_7,x_8\}$ & $\{x_4,x_6,x_9\}$ & $s_1=9,\,s_2=4$ & $x_1,\,x_7$ \\
\bottomrule
\end{tabular}
\end{center}

The three missing elements of $S_r$ (i.e., those with $f_s=0$) are the
$s$-variables substituted as $x_s = -x_{\spec(r)}$ with no additive constant.
For each zero, this is exactly one of the three substitutions.

\subsection{The union: $151$ kills}

Taking the union of the $6\times 55=330$ kills (with overlaps), we find
$151$ distinct coefficients are killed.

\begin{lemma}[Union kill]
\label{lem:union}
The total number of distinct coefficients killed as single-term consequences
of orders $0$ and $2$ of the six zeros is $151$.
\end{lemma}

The explicit list can be obtained by writing out the $35+20=55$ coefficients
per zero and taking their union. Among the $14$ tree triples of
\eqref{eq:trees}, none are killed --- they all survive.

\subsection{Identifying the survivors as tree triples}

We now identify which coefficients are \emph{not} killed by the technique.

\begin{lemma}[Survivors are tree triples]
\label{lem:survivors}
The $14$ coefficients not killed by the procedure of
Section~\ref{sec:proof} are exactly the $14$ tree triples:
\[
a_{135},\,a_{139},\,a_{147},\,a_{149},\,a_{157},\,
a_{246},\,a_{247},\,a_{257},\,a_{258},\,a_{268},\,
a_{358},\,a_{368},\,a_{369},\,a_{469}.
\]
\end{lemma}

\begin{proof}[Proof sketch]
We must show that each of the $151$ non-tree coefficients is killed, and each
of the $14$ tree coefficients is not. The first direction is a case-by-case
check on the union of kill sets across six zeros; we sketch it:

\paragraph{Every non-tree coefficient is killed.}
Suppose $a_{abc}$ is a non-tree coefficient. If the multiset $(a,b,c)$ has a
repeat (double or triple pole), then Lemma~\ref{lem:order0} applied to a
well-chosen $\Z_r$ kills it: one can always choose $r$ such that the repeated
index and the third index all lie in $F_r$. (The $F_r$ sets collectively cover
every index, and each index lies in five of the six $F_r$ sets.)

If $a_{abc}$ has three distinct indices but the corresponding chords are not
a triangulation, then the three chords contain a crossing pair. A careful check
shows that at least one choice of $\Z_r$ places all three indices inside
$F_r\cup S_r\setminus\{x_{\spec(r)}\}$ in exactly the single-$S_r$-index
configuration of Lemma~\ref{lem:order2}.

\paragraph{No tree coefficient is killed.}
A tree coefficient $a_{abc}$ corresponds to a triangulation; its three chords
are a ``fan'' from some hub vertex of the hexagon. A direct case check for each
of the $14$ triangulations shows that none satisfy either the $F_r$ containment
of Lemma~\ref{lem:order0} or the single-$S_r$-index condition of
Lemma~\ref{lem:order2} for any $r\in\{1,\dots,6\}$.

The claim follows. For brevity we omit the full enumeration; it is
straightforward and mechanical.
\end{proof}

\subsection{Concluding the proof of Theorem~\ref{thm:locality}}

Combining Lemmas~\ref{lem:order0}, \ref{lem:order2}, \ref{lem:union}, and
\ref{lem:survivors}, we conclude that on the intersection of the six cyclic
hidden zeros, every non-local coefficient vanishes.

The $14$ tree-triple coefficients are not directly killed by the technique; to
complete the amplitude-uniqueness argument, one must further show that they
are all equal to each other (up to an overall normalization). This comes from
the multi-term equations in the Laurent expansions (orders $1$ and $3$, and
the multi-term parts of orders $0$ and $2$), which we do not tabulate here.

This completes the proof that \emph{locality emerges} from hidden zeros at six
points: every non-local term dies, leaving only triangulations.
\qed

%================================================================
\section{Remarks and pattern toward general $n$}

\subsection{Why the technique is so clean}

The core reason: at each zero $\Z_r$, exactly one free variable
$x_{\spec(r)}$ appears in every substitution. This allows the Laurent
expansion in $x_{\spec(r)}$ to cleanly separate the problem into orders,
each of which is a polynomial identity in only five variables
(the non-special frees).

The leading order (order $0$) performs a \emph{complete kill} of all
coefficients whose denominator doesn't touch the substitution at all. This
is the first $35$ per zero.

The subleading order (order $2$) performs a clean pole-isolation: because
every substituted variable in $S_r$ is a linear function of $x_{\spec(r)}$
(with integer coefficient $\pm 1$), the correction at order $2$ contains an
identifiable new monomial class. This gives $20$ more per zero.

\subsection{Counting the survivors}

The survivors must be coefficients whose indices escape both conditions: they
must have at least one index in $\{x_{\spec(r)}\}$ (or have $\ge 2$ indices in
$S_r$) for every $r$. Working out the overlap over all six zeros, one finds
exactly the $14$ triangulations. A direct proof of this claim --- without
enumeration --- would follow from a combinatorial lemma about the hexagon mesh
structure, which we do not pursue here.

\subsection{Pattern for general $n$}

The structure generalizes. For the $n$-point color-ordered amplitude, the
$n$-gon has $\binom{n}{2}-n$ chords, and there are $n$ cyclic one-zeros, each
constraining three chords in terms of the others. If for each $\Z_r$ one can
identify a special variable whose substitutions are linear in it, the same
Laurent-at-infinity argument should apply, producing a leading-order kill of
all coefficients whose denominator avoids the substitution, and subleading
kills of coefficients with exactly one substituted variable.

At $n=7$, the same Laurent technique should therefore yield a direct proof of
locality, with analogous combinatorics. Verifying this is a natural next task.

\subsection{Relation to existing arguments}

Rodina~\cite{rodina} proved locality and unitarity from hidden zeros using
BCFW shifts and a subset-connectedness induction. The present argument
bypasses both BCFW and induction, using only elementary Laurent-series
manipulation. It is less general in that it applies to the specific ansatz
\eqref{eq:ansatz}, but substantially more elementary where it applies.

\begin{thebibliography}{9}
\bibitem{rodina} B. Rodina, ``Hidden zeros are equivalent to enhanced
ultraviolet scaling and lead to unique amplitudes in $\mathrm{Tr}\,\phi^3$
theory,'' \emph{arXiv:2406.04234}.
\end{thebibliography}

\end{document}
